% Auto-generated from the frozen cohort (do not edit).
\newcommand{\RetfullreplayRow}{full-replay & 63/63 & 4{,}622{,}754 & 0 & 0 & 0 \\}
\newcommand{\RetmaskingRow}{masking & 63/63 & 1{,}994{,}185 & 53{,}583 & 0 & 0 \\}
\newcommand{\RetsummaryresetRow}{summary-reset & 21/63 & 188{,}258 & 0 & 30{,}820 & 280 \\}
\newcommand{\RetboundedhierRow}{bounded-hier & 63/63 & 632{,}699 & 29{,}452 & 0 & 0 \\}
\newcommand{\RetverifiedboardRow}{verified-board & 63/63 & 4{,}977{,}741 & 0 & 0 & 0 \\}
\newcommand{\RetfrontierRow}{frontier & 63/63 & 1{,}906{,}765 & 27{,}728 & 0 & 0 \\}
\newcommand{\RetfrontieradaptRow}{frontier+adapt & 63/63 & 1{,}906{,}765 & 27{,}728 & 0 & 0 \\}
\newcommand{\RetadaptonlyRow}{adapt-only & 63/63 & 1{,}994{,}185 & 53{,}583 & 0 & 0 \\}
\newcommand{\RetfrontierflatRow}{frontier+flat & 63/63 & 1{,}906{,}765 & 27{,}728 & 0 & 0 \\}
\newcommand{\RetfrontierhactRow}{frontier+hact & 63/63 & 1{,}906{,}765 & 27{,}728 & 0 & 0 \\}

% Auto-generated (do not edit).
\newcommand{\FmClusteredRow}{clustered & 312{,}407 & 322{,}810 & 66.8\% & 5{,}696 & 1{,}891 \\}
\newcommand{\FmIndependentRow}{independent & 1{,}047{,}395 & 1{,}120{,}756 & 46.0\% & 47{,}887 & 25{,}837 \\}
\newcommand{\FmGlobalRow}{global & 4{,}622{,}754 & 4{,}276{,}728 & -- & 0 & 0 \\}

% Auto-generated (do not edit).
\newcommand{\EvClusteredRow}{clustered & 58{,}831 & 111{,}031 & +88.7\% \\}
\newcommand{\EvIndependentRow}{independent & 107{,}005 & 159{,}662 & +49.2\% \\}
\newcommand{\EvGlobalRow}{global & 746{,}816 & 96{,}363 & -87.1\% \\}

\section{When Can an Agent Stop Remembering? A Pre-Registered Mechanism Study}\label{sec:retire}

\subsection{The question and why the old framing could not answer it}\label{sec:retire-question}
Version 5.1 of this manuscript answered a serialization question: given
a fixed test-execution schedule, which layout of completion records
crosses a constrained link most cheaply? The companion review's central
objection was that this optimizes an immaterial contributor to agent
cost and cannot, by construction, say anything about when an agent may
safely forget. This section answers the review's stronger question
directly, at the level the evidence can support:

\begin{quote}\itshape
When can an agent stop remembering its execution history without
losing the current evidence and obligations needed to finish correctly,
and what does stopping remember cost and save?
\end{quote}

The study is \emph{mechanism-level}: every number below is produced by
a deterministic engine that models information flow between a context
policy and a seeded work stream. No hosted model is called. The
protocol---arms, cohort, outcomes, and three closure rules---was frozen
in a versioned document before data generation, and all four amendments
were also written before the final cohort ran. Deviations and their
reasons are disclosed in Appendix~\ref{app:retire-deviations}.

\subsection{The frozen protocol}\label{sec:retire-protocol}
\paragraph{Cohort.} Three registry sizes ($n\in\{128,512,1{,}024\}$
component gates), three dependency families reused from the HACT
workload generators (clustered, independent, global), seven seeds
(55011--55017), forty episodes per run: $3\times3\times7=63$ cells.
Every arm consumes the byte-identical episode stream; only the context
policy differs.

\paragraph{Temporal persistence.} The HACT generators draw invalidation
sets i.i.d.\ per episode, but retirement is a temporal question: what
matters is whether a component worked on today is likely to be worked
on again soon. Each episode therefore either persists the working set
with a small shift (probability 0.6) or redraws from the family
generator (probability 0.4). Marginals stay family-matched; the
temporal layer is new and was frozen before data generation
(Amendment~v1.1).

\paragraph{Payloads and footprints.} Each fact carries a payload
(frozen lognormal, median 128\,B) so that history has something to
re-send; each working phase carries a dependency footprint of twelve
gates that persists through persistence steps (the review's exported
interfaces crossing the frontier). Obligations open stochastically on
invalidated gates and come due exactly six episodes later: a policy
that has lost the obligation by then does not act, and the loss is
scored against the policy, never against the stream.

\paragraph{Acceptance contract.} A run is accepted iff no required read
was answered from a stale binding, no obligation was lost before its
due episode, and every due obligation was resolved. Any staleness or
obligation loss rejects the run regardless of byte counts. This is the
frozen quality gate; it is never traded against bytes silently.

\paragraph{Arms.} Ten frozen arms:
(i)~\texttt{full-replay}, the entire history with payloads every call;
(ii)~\texttt{masking}, the observation-masking design: identity-only
context, focused retrieval, no cross-call payload cache;
(iii)~\texttt{summary-reset}, periodic summary and reset every four
episodes, steel-manned as the review demanded: it observes everything
within each period, and its failures are exactly (a) obligations
surviving a reset boundary only through a lossy summary (seeded 50\%
retention per boundary) and (b) window overflow reverting values to
the summary belief, which is silently stale for gates changed after
the last reset;
(iv)~\texttt{bounded-hier}, the vendored Future Code bounded-context
design: a fixed page budget of 64 observations with chronological
retention;
(v)~\texttt{verified-board}, a decentralized status board: one
present-state row per gate;
(vi)~\texttt{frontier}, the proposal: verified-frontier retirement
with a typed capsule (160\,B base, 16\,B per interface binding,
48\,B per pending obligation, 64\,B invariants) whose retained set is
the current working set, its dependency footprint, and pending
obligations;
(vii)~\texttt{frontier+adapt}, frontier retirement plus adaptive
batching (makespan dimension);
(viii)~\texttt{adapt-only}, adaptive batching over masking retention
(attribution ablation);
(ix)~\texttt{frontier+flat}, frontier retirement with a flat evidence
ledger steel-manned as the better of full export and delta export at
64\,B per record;
(x)~\texttt{frontier+hact}, the same with the learned-order
certificate-tree evidence layout (order learned by spectral seriation
on co-invalidation affinity from the first twenty episodes, evaluated
on the held-out second half).

\paragraph{Outcomes and closure rules.} Primary: cost per accepted
project (submitted context bytes plus retrieval bytes over accepted
runs), acceptance rate, staleness defects, obligation losses;
co-primary: retrieval operations, because a live on-demand fetch is a
tool-call round trip whose per-call overhead dominates its payload
bytes. Three closure rules were frozen ex ante:
\emph{Null closure}: if frontier retirement does not reduce
cost-per-accepted by at least 20\% versus the best baseline, the
retirement hypothesis is rejected as a byte-cost claim.
\emph{Attribution closure}: if \texttt{adapt-only} matches
\texttt{frontier+adapt} within 10\%, the contribution is attributed to
scheduling---with the dimension-separated reading mandated ex ante,
since the scheduler is byte-neutral by design.
\emph{HACT closure}: if the tree does not beat the flat steel-man by at
least 5\% on held-out evidence bytes at equal correctness, HACT is
demoted to optional and the paper's center becomes the retirement
policy.

\subsection{Results}\label{sec:retire-results}

\begin{table}[t]\centering\small
\caption{Frozen cohort: 630 runs (63 cells $\times$ 10 arms). Bytes are
per run (submitted context + retrieval); ops are total on-demand
retrieval operations across the 63 cells of each arm; staleness and
loss totals are absolute. Acceptance is the frozen contract: any
staleness or obligation loss rejects the run.}\label{tab:retire-main}
\begin{tabularx}{\linewidth}{@{}l r r r r r@{}}\toprule
Arm & Accepted & Bytes/run & Ops & Stale & Lost\\\midrule
\RetfullreplayRow
\RetmaskingRow
\RetsummaryresetRow
\RetboundedhierRow
\RetverifiedboardRow
\RetfrontierRow
\RetfrontieradaptRow
\RetadaptonlyRow
\RetfrontierflatRow
\RetfrontierhactRow
\bottomrule\end{tabularx}\end{table}

\paragraph{Null closure triggered.} Frontier retirement costs
1{,}906{,}765 bytes per accepted run against full replay's
4{,}622{,}754 (a 58.8\% reduction) but against observation masking's
1{,}994{,}185 it saves only 4.4\% pooled---dominated by the global
family, where masking's context degenerates toward the full registry
and the frontier capsule is comparatively cheap. On the non-global
families the frontier costs \emph{more} bytes than masking
($+3.3\%$ clustered, $+7.0\%$ independent; Table~\ref{tab:retire-fm}).
The pre-registered 20\% gate failed; we therefore claim no byte-cost
productivity win over strong baselines, exactly as the rule requires.
The cheapest arm on submitted bytes, \texttt{summary-reset}, is also
the arm that fails the acceptance contract two-thirds of the time.

\paragraph{What retirement actually buys: correctness and fewer
retrieval round trips.} The falsification result is sharp. Periodic
summary/reset---cheap on bytes---accepts only 21 of 63 runs, with
30{,}820 stale reads and 280 lost obligations: within a reset window
the policy cannot detect that a retained value silently changed, and
obligations crossing a boundary survive only through a lossy summary.
Frontier retirement, masking, the verified board, bounded pages, and
full replay all accept 63/63. Against masking, frontier retirement
cuts on-demand retrieval operations by 48.3\% pooled over non-global
families (66.8\% clustered, 46.0\% independent) at equal correctness.
This is not a byte win---masking's identity-only context is nearly
optimal on bytes---but it is the pre-registered live-cost hypothesis:
in a live cohort, each avoided fetch is an avoided tool-call round
trip, and tool-call overhead (a model invocation) dominates payload
bytes. The live cohort protocol in Section~\ref{sec:live-protocol}
pre-registers this as its primary cost measure.

\paragraph{Attribution closure triggered, with the dimension-separated
reading.} \texttt{adapt-only} matches \texttt{frontier+adapt} on bytes
within 10\%, as expected from a byte-neutral scheduler dimension. Read
dimension-separated, as the protocol mandates: the scheduler affects
makespan only (87.4\% fewer rounds with adaptive batching), and
retirement affects bytes and operations only. Neither dimension
borrows evidence from the other.

\paragraph{HACT survives conditionally, not centrally.}
Table~\ref{tab:retire-evidence} shows the held-out evidence ablation
against the steel-manned flat ledger. Pooled, the learned tree reduces
evidence bytes 59.8\%; per family it \emph{loses} on sparse deltas
($+88.7\%$ clustered, $+49.2\%$ independent) because each packet
carries 128\,B of framing that a delta ledger of a few 64\,B records
undercuts, and it \emph{wins} decisively on dense global updates
($-87.1\%$) because one packet replaces thousands of records. The
tree's 5\% closure gate passes pooled but fails on two of three
families. The honest conclusion is a conditional decision rule, not a
centerpiece: use grouped certificate layouts when update sets are
dense relative to the registry (the global coupling regime, or
whole-epoch re-verification), use flat delta ledgers when updates are
sparse, and select between them by the measured density of the
workload---the same density signal the layout's own training half
already observes. HACT therefore leaves the paper's title and center;
Sections~\ref{sec:compiler}--\ref{sec:eval} retain the component
result with its scope labels unchanged.

\begin{table}[t]\centering\small
\caption{Frontier retirement versus observation masking, per family
(pooled over sizes and seeds). Bytes are per run; the ops column is the
reduction in on-demand retrieval operations. Global coupling leaves no
gap: every read is a current fact for both arms.}\label{tab:retire-fm}
\begin{tabularx}{\linewidth}{@{}l r r r r r@{}}\toprule
Family & Masking B/run & Frontier B/run & Ops reduction & Masking ops & Frontier ops\\\midrule
\FmClusteredRow
\FmIndependentRow
\FmGlobalRow
\bottomrule\end{tabularx}\end{table}

\begin{table}[t]\centering\small
\caption{Held-out evidence-layout ablation (second half of each
stream; order learned on the first half). The flat steel-man is the
better of full and delta export at identical 64\,B record cost.
Positive is a penalty for the tree.}\label{tab:retire-evidence}
\begin{tabularx}{\linewidth}{@{}l r r r@{}}\toprule
Family & Flat steel-man B/run & HACT tree B/run & Difference\\\midrule
\EvClusteredRow
\EvIndependentRow
\EvGlobalRow
\bottomrule\end{tabularx}\end{table}

\subsection{What this study cannot establish}\label{sec:retire-limits}
The engine models information flow, not cognition. A real agent under
observation masking may behave differently than a byte-optimal
retriever; a real summary may retain obligations better or worse than
the seeded 50\% boundary survival; real payloads are not lognormal.
The staleness and obligation-loss results are structural consequences
of each policy's retention mechanism, which is exactly why they were
worth measuring at mechanism level: they are the failure modes the
live cohort must confirm before any productivity claim. Byte numbers
are lower bounds on nothing; they are the frozen accounting of a
frozen model. No hosted model was called; no token or billing claim is
made; the BPE and physical-WAN gaps of Section~\ref{sec:eval} remain
open. The cohort's value is that its three closure rules fired on
pre-registered thresholds, two against our own proposal, and the
manuscript reports the survivors rather than the survivors we hoped
for.
